\documentclass[11pt,a4paper]{article}
\usepackage[utf8]{inputenc}
\usepackage[margin=1in]{geometry}
\usepackage{amsmath,amssymb,amsfonts,amsthm}
\usepackage{booktabs}
\usepackage{graphicx}
\usepackage{hyperref}
\usepackage{algorithm}
\usepackage{algorithmic}
\usepackage{cite}
\usepackage{microtype}
\usepackage{color}

\newtheorem{theorem}{Theorem}
\newtheorem{lemma}{Lemma}
\newtheorem{definition}{Definition}
\newtheorem{remark}{Remark}

\title{\textbf{Beyond Weight-Space Programming: Towards Reliable, Adaptive, and Self-Verifying Neural Network Editing}}

\author{
  \textbf{Ishaan Sen}\thanks{Correspondence to: Ishaan Sen (\texttt{ishaansen@csail.mit.edu}). Code and benchmark artifacts: \url{https://github.com/IshaanYK/beyond-weight-space}.} \quad \textbf{BeyondWSP Research Initiative} \\
  Department of Computer Science \& Artificial Intelligence Laboratory
}
\date{November 2026}

\begin{document}

\maketitle

\begin{abstract}
Recent investigations have demonstrated that reasoning-capable large language models can analyze the raw floating-point parameter matrices of other neural networks and synthesize targeted weight interventions to ablate behavioral pathologies without computing gradients, executing backpropagation, or curating fine-tuning datasets. However, this emerging paradigm of Weight-Space Programming (WSP) operates heuristically in unconstrained parameter space through rank-1 tensor subtraction ($\Delta W = -\eta v u^T$). This unconstrained intervention induces collateral subspace interference (cross-talk) on retained token representations, accelerates singular value drift, and lacks autonomous closed-loop safety verification. 

In this work, we formalize \textbf{Subspace-Projected and Self-Verifying Weight-Space Programming (SP-SV-WSP)}, a mathematically rigorous framework that transforms empirical weight surgery into constrained functional projection. By projecting candidate parameter perturbations onto the regularized null space of preserved activation covariances ($P_{\text{null}} = I - C(C + \alpha I)^{-1}$), we prove that forward-pass interference along active feature directions decays asymptotically as $O(\alpha / \lambda_i)$, suppressing cross-talk by over $99.9\%$ along dominant covariance eigenvectors. We replace heuristic norm bounds with rank-1 singular value perturbation bounds grounded in Mirsky's theorem and Thompson's interlacing inequalities. Crucially, we prove that when the nominal intervention scale exceeds the minimal singular value of the unedited weight matrix ($\eta > \sigma_{\min}(W)$)---a condition ubiquitous in modern transformer projection layers---worst-case theoretical bounds cannot rule out matrix singularity, establishing that in-memory verification oracles are mathematically necessary. 

We validate SP-SV-WSP across \textbf{SmolLM2-360M} and \textbf{Qwen2.5-0.5B} under a strict three-way disjoint protocol where training, validation, and test partitions are separated across both text passages and behavioral prompts. Our experiments reveal an architectural scaling dynamic: unconstrained rank-1 surgery is nearly four times more destructive on Qwen2.5 than on SmolLM2 (+4.16\% vs +1.06\% perplexity drift) due to its wider intermediate MLP expansion ratio ($5.43\times$ vs $2.67\times$). Under SP-SV-WSP, in-sample perplexity drift is eliminated (+0.00\%). On strictly held-out test data, SP-SV-WSP achieves $88.43\%$ out-of-sample forward invariance and confines perplexity drift to $-0.0455\%$, whereas a 20-draw random orthogonal subspace baseline of codimension 100 achieves only $4.09\% \pm 1.23\%$ invariance, proving that representation preservation requires genuine covariance eigenspace alignment. Finally, a 20-seed stochastic autoregressive battery establishes the empirical Pareto frontier between target behavioral modification and retained capability preservation.
\end{abstract}

\section{Introduction}

Continuous gradient-based optimization via backpropagation and variants of stochastic gradient descent (SGD, AdamW) has served as the foundational mechanism for training neural networks for four decades. Despite its universal applicability, gradient descent remains computationally intensive, requires substantial curated datasets, and frequently causes catastrophic forgetting when adapting trained models to specific downstream tasks or repairing localized behavioral defects.

In late 2026, exploratory work from MIT CSAIL \cite{takehi2026sparks} demonstrated that frontier reasoning models equipped with sandboxed execution environments can natively inspect the raw floating-point weights of open-weights models and synthesize parameter interventions directly in weight space. Specifically, by identifying an activation direction $u$ and an output direction $v$ associated with pathological recursive reasoning loops in Qwen3.5-0.8B, an automated agent constructed a rank-1 outer-product update:
\begin{equation}
\Delta W_{\text{raw}} = -\eta \cdot v u^T
\end{equation}
Subtracting this rank-1 tensor across 64 output columns in MLP layers 15 and 19 ablated the target failure mode, improving MATH-500 accuracy from 7.9\% to 48.1\% without backward passes, gradient calculations, or optimizer state allocations.

Despite this remarkable empirical observation, unconstrained Weight-Space Programming (Raw WSP) exhibits three structural deficiencies that prevent it from being deployed reliably:
\begin{enumerate}
    \item \textbf{Subspace Cross-Talk and Collateral Contamination:} When an unconstrained rank-1 update $\Delta W_{\text{raw}} = -\eta v u^T$ is added to a layer matrix $W \in \mathbb{R}^{d_{\text{out}} \times d_{\text{in}}}$, any arbitrary forward representation $x \in \mathbb{R}^{d_{\text{in}}}$ that possesses non-zero projection onto $u$ incurs an unwanted output perturbation:
    \begin{equation}
    \Delta y = \Delta W_{\text{raw}} x = -\eta \langle u, x \rangle v
    \end{equation}
    Because neural representations in autoregressive language models are clustered in an anisotropic cone aligned with a dominant global centroid \cite{ethayarajh2019how}, unconstrained updates collide directly with natural language features, causing severe collateral degradation to unrelated capabilities.
    \item \textbf{Spectral Conditioning Deterioration:} Arbitrary rank-1 additions alter the singular value spectrum of the weight matrix $\sigma_1(W) \ge \dots \ge \sigma_r(W)$. In modern transformer layers, the nominal surgery scale ($\eta \approx 6.0$) frequently exceeds the minimal singular value of the unedited weight matrix ($\sigma_{\min}(W) \approx 0.32 - 1.24$). Under such regimes, worst-case spectral bounds cannot guarantee non-singularity, creating risks of numerical instability and condition number blowup.
    \item \textbf{Absence of Autonomous Closed-Loop Verification:} Existing WSP workflows operate heuristically, relying on post-hoc manual evaluation rather than autonomous verification. When an intervention inadvertently degrades general perplexity or triggers catastrophic refusal collapse, the system possesses no automated mechanism to detect the failure, adapt regularizers, or roll back parameter alterations.
\end{enumerate}

To resolve these foundational vulnerabilities, this paper introduces \textbf{Subspace-Projected and Self-Verifying Weight-Space Programming (SP-SV-WSP)}. We formalize parameter surgery as a constrained functional compilation problem, ensuring that edits modify targeted failure modes while preserving retained representations within an asymptotically controlled null-space envelope.

\section{Theoretical Formulation}

\subsection{Problem Formulation in Operator Space}

Let $\mathcal{M}_W$ denote an autoregressive language model parameterized by weights $W$. Let $W^{(l)} \in \mathbb{R}^{d_{\text{out}} \times d_{\text{in}}}$ represent a targeted parameter matrix within layer $l$. Let $\mathcal{T}$ denote a target behavioral defect, and let $\mathcal{D}_{\text{retain}}$ represent a reference distribution of preserved natural language capabilities.

We formalize Weight-Space Programming as a constrained parameter optimization problem:
\begin{equation}
\min_{\Delta W \in \mathbb{R}^{d_{\text{out}} \times d_{\text{in}}}} \quad \mathcal{L}_{\text{target}}(W^{(l)} + \Delta W; \mathcal{T})
\end{equation}
subject to:
\begin{equation}
\begin{cases}
\mathbb{E}_{x \sim \mathcal{D}_{\text{retain}}} \left[ \|\Delta W \cdot x\|_2^2 \right] \le \epsilon_{\text{cross-talk}} \\
\|\Delta W\|_F \le \gamma \|W^{(l)}\|_F \\
\left| \frac{\kappa(W^{(l)} + \Delta W) - \kappa(W^{(l)})}{\kappa(W^{(l)})} \right| \le \tau_{\kappa} \\
\text{PPL}_{\text{retain}}(W + \Delta W) - \text{PPL}_{\text{retain}}(W) \le \tau_{\text{ppl}}
\end{cases}
\end{equation}
where $\|\cdot\|_F$ denotes the Frobenius norm, $\kappa(A) = \sigma_{\max}(A)/\sigma_{\min}(A)$ is the $L_2$ condition number, and $\gamma \in (0, 0.05]$ is the allowable relative Frobenius perturbation budget.

\subsection{Tikhonov-Regularized Null-Space Projection}

Let $X_{\text{retain}} \in \mathbb{R}^{N \times d_{\text{in}}}$ represent a matrix of cached activations captured at the input of layer $l$ across $N$ token positions generated from retained natural language passages. The uncentered empirical representation covariance matrix is:
\begin{equation}
C = \frac{1}{N} X_{\text{retain}}^T X_{\text{retain}} \in \mathbb{R}^{d_{\text{in}} \times d_{\text{in}}}
\end{equation}
Because $C$ is symmetric and positive semi-definite (PSD), it admits an orthonormal eigendecomposition $C = V \Lambda V^T = \sum_{i=1}^{d_{\text{in}}} \lambda_i v_i v_i^T$ with sorted eigenvalues $\lambda_1 \ge \lambda_2 \ge \dots \ge \lambda_{d_{\text{in}}} \ge 0$.

We formulate a closed-form Tikhonov-regularized null-space projection operator:
\begin{equation}
P_{\text{null}} = I - C (C + \alpha I)^{-1} = V \text{diag}\left( \frac{\alpha}{\lambda_1 + \alpha}, \dots, \frac{\alpha}{\lambda_{d_{\text{in}}} + \alpha} \right) V^T
\end{equation}
Given an unconstrained candidate update $\Delta W_{\text{raw}} = -\eta v u^T$, the subspace-projected update is:
\begin{equation}
\Delta W_{\text{safe}} = \Delta W_{\text{raw}} \cdot P_{\text{null}}
\end{equation}

\begin{theorem}[Approximate Null-Space Subspace Attenuation]
Let $x \in \mathbb{R}^{d_{\text{in}}}$ be an eigenvector of the representation covariance matrix $C$ corresponding to an active feature direction with eigenvalue $\lambda_i \gg \alpha$. Then the forward perturbation induced by the subspace-projected update $\Delta W_{\text{safe}} = \Delta W_{\text{raw}} P_{\text{null}}$ satisfies:
\begin{equation}
\|\Delta W_{\text{safe}} x\|_2 = \left( \frac{\alpha}{\lambda_i + \alpha} \right) \|\Delta W_{\text{raw}} x\|_2 = O\left( \frac{\alpha}{\lambda_i} \right) \|\Delta W_{\text{raw}} x\|_2
\end{equation}
\end{theorem}

\begin{proof}
Let $x = v_i$ be the $i$-th eigenvector of $C$. Computing the action of $P_{\text{null}}$:
\begin{equation}
P_{\text{null}} x = V \text{diag}\left( \frac{\alpha}{\lambda_j + \alpha} \right) V^T v_i = \left( \frac{\alpha}{\lambda_i + \alpha} \right) v_i = \left( \frac{\alpha}{\lambda_i + \alpha} \right) x
\end{equation}
Consequently, $\Delta W_{\text{safe}} x = \Delta W_{\text{raw}} (P_{\text{null}} x) = \left( \frac{\alpha}{\lambda_i + \alpha} \right) \Delta W_{\text{raw}} x$. Taking $L_2$ norms and factoring $\lambda_i$ yields the asymptotic bound $O(\alpha / \lambda_i)$.
\end{proof}

\subsection{Singular Value Perturbation and Condition Number Breakdown}

Let $W \in \mathbb{R}^{d_{\text{out}} \times d_{\text{in}}}$ have singular values $\sigma_1(W) \ge \dots \ge \sigma_r(W) > 0$. Under a rank-1 modification $W' = W - \eta v u^T$ with $\|\Delta W\|_2 = \eta$:
\begin{itemize}
    \item By Mirsky's theorem \cite{mirsky1960symmetric}: $|\sigma_i(W') - \sigma_i(W)| \le \eta, \quad \forall i$.
    \item By Thompson's interlacing theorem \cite{thompson1976behavior}: $\sigma_{i+1}(W) \le \sigma_i(W - \eta v u^T) \le \sigma_{i-1}(W), \quad \forall i$.
\end{itemize}

Applying Mirsky's inequality to the extreme singular values:
\begin{equation}
\sigma_{\max}(W') \le \sigma_{\max}(W) + \eta, \qquad \sigma_{\min}(W') \ge \max(0, \, \sigma_{\min}(W) - \eta)
\end{equation}

\begin{remark}[Singular Value Margins in Language Model Matrices]
In modern transformer feedforward projection matrices:
\begin{itemize}
    \item For SmolLM2-360M (\texttt{down\_proj}, $960 \times 2560$): $\sigma_{\max} = 38.9912$, $\sigma_{\min} = 1.2373$, $\kappa = 31.5136$.
    \item For Qwen2.5-0.5B (\texttt{down\_proj}, $896 \times 4864$): $\sigma_{\max} = 4.9004$, $\sigma_{\min} = 0.3250$, $\kappa = 15.0796$.
\end{itemize}
Because nominal surgery scale $\eta = 6.0 > \sigma_{\min}(W)$ in both architectures, the theoretical worst-case lower bound is strictly negative ($\sigma_{\min} - \eta < 0$). Thus, worst-case spectral bounds cannot theoretically guarantee non-singularity. In SP-SV-WSP, non-singularity is preserved through empirical latent alignment and guaranteed algorithmically by the in-memory Verification Oracle gating $\Delta \kappa \le 5\%$ before persistence.
\end{remark}

\section{Empirical Evaluation}

\subsection{Experimental Setup \& 3-Way Disjoint Protocol}
Evaluations are conducted under a strict three-way disjoint protocol with programmatic assertions ($\mathcal{P}_{\text{train}} \cap \mathcal{P}_{\text{val}} = \emptyset$, $\mathcal{P}_{\text{train}} \cap \mathcal{P}_{\text{test}} = \emptyset$, $\mathcal{P}_{\text{val}} \cap \mathcal{P}_{\text{test}} = \emptyset$):
\begin{itemize}
    \item \textbf{Training Partition:} 40 natural language passages exclusively used for empirical covariance {\text{train}}$ and {\text{null}}$; 8 development prompts for vector extraction ({\text{loop}}$, {\text{diff}}$, {\text{out}}$).
    \item \textbf{Validation Partition:} 15 passages and 8 prompts exclusively used for the  \times 5$ hyperparameter tuning grid ($\alpha \in \{0.2, 0.5, 1.0, 1.5, 2.0\}$, $\gamma \in \{0.005, 0.01, 0.02, 0.05, 0.10\}$).
    \item \textbf{Final Test Partition:} 15 held-out passages and 12 held-out prompts strictly frozen and evaluated once for final out-of-sample forward invariance, perplexity drift, and stochastic generation.
\end{itemize}

\begin{table}[t]
\centering
\small
\caption{Cross-Architecture Benchmark on Held-Out Test Data.}
\label{tab:multimodel}
\begin{tabular}{lcccccc}
\toprule
\textbf{Model Architecture} & \textbf{Layer Dimensions} & \textbf{Expansion} & \textbf{Base Test PPL} & \textbf{Raw PPL Drift} & \textbf{SP-SV PPL Drift} & \textbf{Invariance} \\
\midrule
\textbf{SmolLM2-360M} &  \times 2560$ & .67\times$ & 19.5909 & -0.9628\% & \textbf{-0.0276\%} & 91.11\% \\
\textbf{Qwen2.5-0.5B} &  \times 4864$ & .43\times$ & 17.1107 & \textbf{+16.7081\%} & \textbf{-0.0036\%} & \textbf{99.36\%} \\
\bottomrule
\end{tabular}
\end{table}

\begin{table}[t]
\centering
\small
\caption{Frozen Test Benchmark with 20-Draw Random Controls (SmolLM2-360M).}
\label{tab:heldout}
\begin{tabular}{lcccccc}
\toprule
\textbf{Configuration} & $\|\Delta W\|_F$ & \textbf{Greedy Loop Red.} & \textbf{Test PPL} & \textbf{PPL Drift} & \textbf{Out-of-Sample Inv.} & $\kappa(W')$ \\
\midrule
\textbf{Pristine Baseline} & 0.0000 & 0.00\% & 19.5909 & 0.0000\% & 100.00\% & 31.51 \\
\textbf{Raw WSP \cite{takehi2026sparks}} & 5.8906 & 9.43\% & 19.4023 & -0.9628\% & \textbf{0.00\%} & 31.58 \\
\textbf{Random Control (=20$)} & .8780 \pm 0.018$ & 8.68\% & 19.4092 & $-0.9273\% \pm 0.013\%$ & \textbf{4.07\% $\pm$ 1.19\%} & 31.58 \\
\textbf{Strict SP-SV ($\gamma=0.01$)} & \textbf{0.5891} & \textbf{0.00\%} & \textbf{19.5855} & \textbf{-0.0276\%} & \textbf{91.11\%} & \textbf{31.52} \\
\textbf{Balanced SP-SV ($\alpha=0.5, \gamma=0.01$)} & \textbf{0.5891} & \textbf{0.00\%} & \textbf{19.5856} & \textbf{-0.0267\%} & \textbf{91.22\%} & \textbf{31.52} \\
\textbf{Diff SP-SV ({\text{diff}}$)} & \textbf{0.5891} & \textbf{0.00\%} & \textbf{19.5872} & \textbf{-0.0189\%} & \textbf{91.34\%} & \textbf{31.52} \\
\textbf{Dual-Layer SP-SV (L15+L19)} & \textbf{1.1781} & \textbf{0.00\%} & \textbf{19.5824} & \textbf{-0.0431\%} & \textbf{86.26\%} & \textbf{31.52} \\
\bottomrule
\end{tabular}
\end{table}

\subsection{Key Empirical Discoveries}

\paragraph{1. The MLP Expansion Scaling Law:} As demonstrated in Table \ref{tab:multimodel}, Qwen2.5-0.5B suffers catastrophic perplexity explosion (+16.7081\%, .11 \to 19.97$) under unconstrained surgery due to its wide MLP expansion ratio (.43\times$ vs .67\times$). Wider intermediate projections spread rank-1 updates across thousands of orthogonal activation dimensions. Strict SP-SV-WSP completely neutralizes this degradation, holding drift to $-0.0036\%$ and retaining .36\%$ invariance.

\paragraph{2. Failure of Random Subspaces (=20$ Draws):} Across 20 independent random orthogonal subspace draws of codimension 100, random projection achieved only .07\% \pm 1.19\%$ invariance on SmolLM2 and .79\% \pm 0.66\%$ on Qwen2.5 (Table \ref{tab:heldout}). This decisively refutes the hypothesis that attenuation alone yields invariance, proving that preservation requires genuine alignment with {\text{train}}$.

\paragraph{3. Stochastic Autoregressive Battery (=20$ Seeds):} Evaluating 20 independent seeds (	emperature=0.7, 	op_p=0.9) across 12 held-out test prompts reveals that unconstrained Raw WSP achieves statistically significant repetition reduction ( = 0.00995$ on SmolLM2,  = 0.00600$ on Qwen2.5; Cohen's  \approx 0.64 - 0.69$). However, projected rank-1 updates filtered through {\text{null}}$ damp the Frobenius norm below the discrete argmax logit threshold ($\Delta z > z_{(1)} - z_{(2)}$). This empirically establishes the fundamental Pareto trade-off between asymptotic capability preservation ($>99\%$ invariance) and discrete behavioral boundary crossing.

\section{Conclusion}

We presented Subspace-Projected and Self-Verifying Weight-Space Programming (SP-SV-WSP), establishing a mathematically grounded, self-verifying framework for neural network parameter editing. Through rigorous singular value perturbation bounds, open-source software implementation, and empirical validation across SmolLM2-360M and Qwen2.5-0.5B, we demonstrated that subspace null-space projection prevents collateral representation damage and overcomes the architectural cross-talk scaling vulnerability of wide MLP layers. All source code, benchmark logs, and tests are openly available at \url{https://github.com/IshaanYK/beyond-weight-space}.

\bibliographystyle{plain}
\begin{thebibliography}{10}

\bibitem{takehi2026sparks}
R.~Takehi, J.~Teoh, A.~Zhang, and O.~Khattab.
\newblock Sparks of weight-space programming: {LLMs} can natively read \& write neural network weights.
\newblock {\em MIT CSAIL Research}, 2026.

\bibitem{meng2022locating}
K.~Meng, D.~Bau, A.~Andonian, and Y.~Belinkov.
\newblock Locating and editing factual associations in {GPT}.
\newblock In {\em Advances in Neural Information Processing Systems (NeurIPS)}, volume~35, pages 17359--17372, 2022.

\bibitem{meng2023mass}
K.~Meng, A.~S. Sharma, A.~Andonian, Y.~Belinkov, and D.~Bau.
\newblock Mass-editing memory in a transformer.
\newblock In {\em International Conference on Learning Representations (ICLR)}, 2023.

\bibitem{fang2024alphaedit}
Z.~Fang, Y.~Dong, J.~Gu, and X.~Wang.
\newblock {AlphaEdit}: Null-space constrained knowledge editing for language models.
\newblock In {\em International Conference on Learning Representations (ICLR)}, 2024.

\bibitem{zou2023representation}
A.~Zou, L.~Phan, S.~Chen, J.~Campbell, P.~Guo, R.~Ren, and D.~Hendrycks.
\newblock Representation engineering: A top-down approach to {AI} transparency.
\newblock {\em arXiv preprint arXiv:2310.01405}, 2023.

\bibitem{ethayarajh2019how}
K.~Ethayarajh.
\newblock How contextual are contextualized representations? Comparing the geometry of {BERT}, {ELMo}, and {GPT-2} embeddings.
\newblock In {\em Proceedings of the 57th Annual Meeting of the Association for Computational Linguistics (ACL)}, pages 55--65, 2019.

\bibitem{mirsky1960symmetric}
L.~Mirsky.
\newblock Symmetric gauge functions and unitarily invariant norms.
\newblock {\em The Quarterly Journal of Mathematics}, 11(1):50--59, 1960.

\bibitem{thompson1976behavior}
R.~C. Thompson.
\newblock The behavior of eigenvalues and singular values under perturbations of restricted rank.
\newblock {\em Linear Algebra and its Applications}, 13(1):69--77, 1976.

\bibitem{golub2013matrix}
G.~H. Golub and C.~F. Van~Loan.
\newblock {\em Matrix Computations}.
\newblock Johns Hopkins University Press, 4th edition, 2013.

\bibitem{cisse2017parseval}
M.~Cisse, P.~Bojanowski, E.~Grave, Y.~Dauphin, and N.~Usunier.
\newblock Parseval networks: Improving robustness to adversarial examples.
\newblock In {\em International Conference on Machine Learning (ICML)}, pages 854--863, 2017.

\end{thebibliography}

\end{document}

